\subsubsection{Barra pelota}
Un brazo de longitud \( l \) está articulado en uno de sus extremos al centro de un disco circular de radio \( r \), como se muestra en la Figura \ref{fig:rodadura_varilla}. El disco rueda sin deslizamiento a lo largo de una superficie horizontal fija. Usando la variable \( x \) para describir la posición del centro del disco y la variable \( \phi \) para describir la orientación del brazo con respecto a la vertical, determine los siguientes vectores según lo visto por un observador fijo al suelo:

\begin{enumerate}
    \item La velocidad y aceleración del centro del disco.
    \item La velocidad y aceleración angular del disco
    \item La velocidad y aceleración del extremo libre del brazo.
\end{enumerate}
\begin{figure}[!ht]
    \centering
    \begin{tikzpicture}[>=Latex, scale=1.8]
        % --- Parámetros Geométricos ---
        \def\R{1.5}      % Radio del cilindro
        \def\L{2.2}      % Longitud de la varilla l
        \def\PhiAng{30}  % Ángulo phi respecto a la vertical

        % --- 1. Suelo (Superficie de rodadura) ---
        \draw[thick] (-2,0) -- (3.5,0);
        \fill[pattern=north east lines, pattern color=gray!60] (-2, -0.2) rectangle (3.5, 0);

        % --- 2. Cilindro ---
        \coordinate (O) at (0, \R); % Centro O
        \coordinate (Q) at (0, 0);  % Punto de contacto Q
        
        % Sombreado degradado para el cilindro
        \shadedraw[inner color=gray!10, outer color=gray!70] (O) circle (\R);
        \filldraw (O) circle (1.5pt) node[above left] {$O$};
        \filldraw (Q) circle (1pt) node[above left, xshift=-2pt] {$Q$};

        % --- 3. Ejes y Referencias ---
        % Eje x desde el centro
        \draw[->, thick] (O) -- (2.5, \R) node[below] {$x$};
        % Línea vertical discontinua
        \draw[dashed, gray] (O) -- (0, -0.5);
        % Radio R
        \draw[->, thick] (O) -- ({\R*cos(45)}, {\R+\R*sin(45)}) node[midway, above left] {$R$};

        % --- 4. Varilla acoplada en O ---
        % Punto P al final de la varilla
        \coordinate (P) at ({\L*sin(\PhiAng)}, {\R - \L*cos(\PhiAng)});
        
        % Varilla con grosor (estilo 2D)
        \draw[line width=6pt, line cap=round, color=black] (O) -- (P);
        \draw[line width=5pt, line cap=round, color=gray!50!white] (O) -- (P);
        \draw[thick] (O) -- (P);
        \filldraw (P) circle (1pt) node[below right] {$P$};
        \node[right, xshift=3pt] at ($(O)!0.5!(P)$) {$l$};

        % --- 5. Ángulo Phi ---
        \draw[->] (0, 0.5) arc (-90:{-90+\PhiAng}:0.8);
        \node at (0.3, 0.4) {$\phi$};

        % --- 6. Indicación "No desliza" ---
        \coordinate (NoSlipPoint) at (-0.8, 0.2);
        \node[left] at (-1, 0.6) {No desliza};
        \draw[thin] (-1, 0.5) to[out=0, in=120] (-0.3, 0.05);

    \end{tikzpicture}
    \caption{Cilindro de radio $R$ rodando sin deslizar sobre una superficie horizontal. Una varilla de longitud $l$ está articulada en el centro $O$, formando un ángulo $\phi$ con la vertical.}
    \label{fig:rodadura_varilla}
\end{figure}